\documentclass[12pt]{article}
\usepackage[utf8]{inputenc}
\usepackage[spanish]{babel}
\usepackage{amsmath}
\usepackage{amsfonts}
\usepackage{amssymb}
\usepackage{geometry}
\geometry{letterpaper, margin=2.5cm}

\begin{document}

\begin{center}
    \Large \textbf{Evaluación de Examen 1: Unidad I} \\
    \large Física para Ciencias de la Salud (005-1713) \\
    \vspace{0.5cm}
    \textbf{Estudiante:} Shantal Romero \quad \textbf{C.I.:} 33.093.620
    \vspace{0.3cm} \\
    \large \textbf{Calificación Final: 9/10 puntos}
\end{center}

\vspace{0.5cm}

\section*{Análisis de Resultados}

\subsection*{1. Ángulo plano (Conversión a radianes)}
\textbf{Procedimiento y resultado:} Plantea correctamente la conversión empleando el factor equivalente de $180^\circ \rightarrow \pi$ rad. Desarrolla la simplificación de la fracción $\left(\frac{75\pi}{180} = \frac{5\pi}{12}\right)$ y adicionalmente calcula su aproximación decimal ($0,42\pi$ rad). Procedimiento impecable. \\
\textbf{Veredicto:} \textbf{Correcto} (2/2 pts).

\subsection*{2. Sistemas de coordenadas (Longitud del hueso)}
\textbf{Procedimiento y resultado:} Utiliza la fórmula formal de la distancia euclidiana. Al sustituir los valores del eje $y$, comete un levísimo error de orden escribiendo $(2-10)^2$ en lugar de $(10-2)^2$. Sin embargo, este detalle no afecta el resultado matemático, ya que al elevar al cuadrado obtiene correctamente $64$. Finaliza el cálculo sin problemas ($\sqrt{36+64} = \sqrt{100}$) llegando a $10$ cm. \\
\textbf{Veredicto:} \textbf{Correcto} (2/2 pts).

\subsection*{3. Vectores (Fuerza resultante)}
\textbf{Procedimiento y resultado:} La estudiante demuestra mucho orden. Plantea la suma de ambos vectores ($\vec{F}_R = \vec{F}_1 + \vec{F}_2$) y agrupa las componentes $\hat{i}$ ($45-15$) y $\hat{j}$ ($25+35$) por separado. El resultado final ensamblado es matemáticamente perfecto y está acompañado de su unidad respectiva: $(30\hat{i} + 60\hat{j})$ N. \\
\textbf{Veredicto:} \textbf{Correcto} (2/2 pts).

\subsection*{4. Potencias de diez (Notación científica y prefijo)}
\textbf{Procedimiento y resultado:} Transforma la cifra decimal correctamente utilizando potencias de base diez ($12,5 \times 10^{-6}$ m). No obstante, como ha sido la tendencia generalizada en el grupo, omite la instrucción adicional de la pregunta que solicitaba identificar el nombre explícito del prefijo del S.I. (micrómetros). \\
\textbf{Veredicto:} \textbf{Incompleto / Parcialmente correcto} (1/2 pts).

\subsection*{5. Sistemas de unidades (Conversión de densidad)}
\textbf{Procedimiento y resultado:} Plantea los factores de conversión en cadena de forma excelente, escribiendo a un lado las equivalencias teóricas exactas ($1 \text{ kg} = 1000 \text{ g}$ y $1 \text{ m}^3 = 1.000.000 \text{ cm}^3$). Deduce correctamente que tras simplificar los ceros solo hace falta multiplicar por 1000. Su cálculo final arroja la cifra correcta de $1030 \text{ Kg/m}^3$. \\
\textbf{Veredicto:} \textbf{Correcto} (2/2 pts).

\vspace{0.5cm}
\section*{Comentarios Generales y Sugerencias}
¡Muy buen examen! La estudiante muestra un nivel analítico muy sólido, manteniendo un gran orden en la presentación de sus fórmulas y resultados.
\begin{itemize}
    \item Se recomienda cuidar minuciosamente el orden de las coordenadas ($y_2 - y_1$) en la fórmula de distancia para evitar inconsistencias de notación, a pesar de que el exponente par enmascare el error de signo (Pregunta 2).
    \item Recordar siempre leer hasta el final los enunciados para no omitir respuestas teóricas (como el nombre del prefijo en la Pregunta 4).
\end{itemize}

\end{document}